% from ega1-7-fr.tex
\oldpage[I]{160}
Il faut prouver que $p\circ\overline f$ est l'identité dans le domaine de
définition de $\overline f$~; puisque $X$ est séparé sur $S$, cela résulte
encore de (\hyperref[I.7.2.2.1-fr]{7.2.2.1}).

\begin{corollary}[7.2.6]
\label{I.7.2.6-fr}
Soient $X$ un préschéma réduit, $U$ un ouvert partout dense de $X$. Il y a
correspondance biunivoque canonique entre les sections de $\mathscr{O}_X$ au-dessus
de $U$ et les fonctions rationnelles $f$ sur $X$ définies en tout point de $U$.
\end{corollary}

Compte tenu de (\hyperref[I.7.2.3-fr]{7.2.3}),
(\hyperref[I.7.1.2-fr]{7.1.2}) et (\hyperref[I.7.1.3-fr]{7.1.3}), il suffit de
remarquer que le $X$-préschéma $X\otimes_{\mathbf Z}\mathbf Z[T]$ est séparé
au-dessus de $X$ (\hyperref[I.5.5.1-fr]{5.5.1, (iv)}).

\begin{corollary}[7.2.7]
\label{I.7.2.7-fr}
Soient $Y$ un préschéma réduit, $f:X\to Y$ un morphisme séparé, $U$ un ouvert
partout dense de $Y$, $g:U\to f^{-1}(U)$ une $U$-section de $f^{-1}(U)$, $Z$
le sous-préschéma réduit de $X$ ayant $\overline{g(U)}$ pour espace sous-jacent
(\hyperref[I.5.2.1-fr]{5.2.1}). Pour que $g$ soit restriction d'une
$Y$-section de $X$ (autrement dit (\hyperref[I.7.2.5-fr]{7.2.5}) pour que
l'application rationnelle de $Y$ dans $X$ prolongeant $g$ soit partout
définie), il faut et il suffit que la restriction de $f$ à $Z$ soit un
isomorphisme de $Z$ sur $Y$.
\end{corollary}

La restriction de $f$ à $f^{-1}(U)$ est un morphisme séparé
(\hyperref[I.5.5.1-fr]{5.5.1, (i)}), donc $g$ est une immersion fermée
(\hyperref[I.5.4.6-fr]{5.4.6}), et par suite
$g(U)=Z\cap f^{-1}(U)$ et le sous-préschéma induit par $Z$ sur l'ouvert
$g(U)$ de $Z$ est identique au sous-préschéma fermé de $f^{-1}(U)$ associé à
$g$ (\hyperref[I.5.2.1-fr]{5.2.1}). Il est clair alors que la condition de
l'énoncé est suffisante, car si elle est remplie et si $f_Z:Z\to Y$ est la
restriction de $f$ à $Z$ et $\overline g:Y\to Z$ l'isomorphisme réciproque,
$\overline g$ prolonge $g$. Inversement, si $g$ est restriction à $U$ d'une
$Y$-section $h$ de $X$, $h$ est une immersion fermée
(\hyperref[I.5.4.6-fr]{5.4.6}), donc $h(Y)$ est fermé, et comme il est contenu
dans $Z$, il est égal à $Z$, et il résulte de
(\hyperref[I.5.2.1-fr]{5.2.1}) que $h$ est nécessairement un isomorphisme de
$Y$ sur le sous-préschéma fermé $Z$ de $X$.

\begin{env}[7.2.8]
\label{I.7.2.8-fr}
Soient $X$, $Y$ deux $S$-préschémas, $X$ étant supposé réduit et $Y$ séparé
sur $S$. Soit $f$ une $S$-application rationnelle de $X$ dans $Y$, et soit
$x$ un point de $X$~; on peut composer $f$ avec le $S$-morphisme canonique
$\operatorname{Spec}(\mathscr{O}_x)\to X$
(\hyperref[I.2.4.1-fr]{2.4.1}) pourvu que la trace sur
$\operatorname{Spec}(\mathscr{O}_x)$ du domaine de définition de $f$ soit
dense dans $\operatorname{Spec}(\mathscr{O}_x)$ (identifié à l'ensemble des
$z\in X$ tels que $x\in\overline{\{z\}}$
(\hyperref[I.2.4.2-fr]{2.4.2})). Ceci aura lieu dans les cas suivants~:
\begin{enumerate}
  \item[1\textsuperscript{o}] $X$ est irréductible (donc intègre), car alors
    le point générique $\xi$ de $X$ est le point générique de
    $\operatorname{Spec}(\mathscr{O}_x)$~; comme le domaine de définition $U$
    de $f$ contient $\xi$, $U\cap\operatorname{Spec}(\mathscr{O}_x)$ contient
    $\xi$, donc est dense dans $\operatorname{Spec}(\mathscr{O}_x)$.
  \item[2\textsuperscript{o}] $X$ est localement noethérien~; notre assertion
    résulte en effet alors du
\end{enumerate}
\end{env}

\begin{lemma}[7.2.8.1]
\label{I.7.2.8.1-fr}
Soient $X$ un préschéma dont l'espace sous-jacent est localement noethérien,
$x$ un point de $X$. Les composantes irréductibles de
$\operatorname{Spec}(\mathscr{O}_x)$ sont les traces sur
$\operatorname{Spec}(\mathscr{O}_x)$ des composantes irréductibles de $X$
contenant $x$. Pour qu'un ouvert $U\subset X$ soit tel que
$U\cap\operatorname{Spec}(\mathscr{O}_x)$ soit dense dans
$\operatorname{Spec}(\mathscr{O}_x)$, il faut et il suffit qu'il rencontre les
composantes irréductibles de $X$ contenant $x$ (ce qui a lieu en particulier
si $U$ est dense dans $X$).
\end{lemma}

La seconde assertion résulte évidemment de la première, et il suffit donc de
démontrer celle-ci. Comme $\operatorname{Spec}(\mathscr{O}_x)$ est contenu
dans tout ouvert affine $U$ contenant $x$, et que les composantes
irréductibles de $U$ contenant $x$ sont les traces sur $U$ des composantes
irréductibles de $X$ contenant $x$
(\hyperref[0.2.1.6-fr]{0, 2.1.6}), on peut supposer $X$ affine d'anneau $A$.
Comme les idéaux premiers de $A_x$ correspondent biunivoquement aux idéaux
premiers
